\documentclass{article}



\usepackage{arxiv}

\usepackage[utf8]{inputenc} % allow utf-8 input
\usepackage[T1]{fontenc}    % use 8-bit T1 fonts
\usepackage{hyperref}       % hyperlinks
\usepackage{url}            % simple URL typesetting
\usepackage{booktabs}       % professional-quality tables
\usepackage{amsfonts}       % blackboard math symbols
\usepackage{nicefrac}       % compact symbols for 1/2, etc.
\usepackage{microtype}      % microtypography
\usepackage{lipsum}		% Can be removed after putting your text content
\usepackage{graphicx}
\usepackage{natbib}
\usepackage{doi}

\usepackage{xcolor}
\hypersetup{
    colorlinks,
    linkcolor={red!50!black},
    citecolor={blue!50!black},
    urlcolor={blue!80!black}
}

%===== Aqui tienes que escribir el título del trabajo =======%
\title{Esto es un Titulo muy bonito de la investigacion}

%=====  Aqui debes poner la Fecha   =======%
\date{Diciembre, 2023}	% Aqui puede cambiar la fecha que aparece en el pdf.



%===== Aqui van los nombres de lo miembros del grupo de investigacion    =======%
\author{ 
{\includegraphics[scale=0.2]{star.png}\hspace{1mm}
        Yerko Novoa	                     }\\
	Poder de las Mareas Cosmicas          \\
	Sabado                                \\
	PentaUC                               \\
	\texttt{Correo@pentauc.cl}            \\
\And
{\includegraphics[scale=0.2]{star.png}\hspace{1mm}
        Roronoa Zoro	                 	 }\\
	Poder de las Mareas Cosmicas          \\
	Sabado                                \\
	PentaUC                               \\
	\texttt{Correo@pentauc.cl}            \\
\And
{\includegraphics[scale=0.2]{star.png}\hspace{1mm}
        Gol D. Roger		             	 }\\
 	Poder de las Mareas Cosmicas          \\
	Sabado                                \\
	PentaUC                               \\
	\texttt{Correo@pentauc.cl}            \\
\And
{\includegraphics[scale=0.2]{star.png}\hspace{1mm}
        Gol Big Mom	Charlotte	         	 }\\
	Poder de las Mareas Cosmicas          \\
	Sabado                                \\
	PentaUC                               \\
	\texttt{Correo@pentauc.cl}            \\
\And
{\includegraphics[scale=0.2]{star.png}\hspace{1mm}
        Sanji	Vinsmoke	                 }\\
	Poder de las Mareas Cosmicas          \\
	Sabado                                \\
	PentaUC                               \\
	\texttt{Correo@pentauc.cl}            \\
}


%=====  Esto no se toca   =======%
% Uncomment to override  the `A preprint' in the header
%\renewcommand{\headeright}{Postprint}
%\renewcommand{\undertitle}{Postprint}
\renewcommand{\shorttitle}{Penta Review Letters UC}

\begin{document}
\maketitle

%=====  Abstract significa resumen: Esto es un resumen de todo el trabajo   =======%
\begin{abstract}

     aca va el resumen de nuestro trabajo
    Suelen tener entre 200-300 palabras, y debe ser lo suficientemente claro y conciso como para que un lector interesado pueda comprender rápidamente el tema del artículo y sus principales conclusiones.\\
    Las características de un abstract de un paper de física teórica son las siguientes: \\
    \begin{itemize}
        \item Concisión: El abstract debe ser conciso y directo. Debe proporcionar una descripción clara y concisa del contenido del artículo.
        \item Precisión: El abstract debe ser preciso y completo. Debe proporcionar una descripción precisa del contenido del artículo, sin omitir información importante.
        \item Objetividad: El abstract debe ser objetivo. Debe evitar opiniones personales o juicios de valor.
        \item Claridad: El abstract debe ser claro y fácil de entender. Debe utilizar un lenguaje sencillo y evitar jerga técnica innecesaria.
    \end{itemize}
\end{abstract}


%=====  Esto es para poner palabras claves del Resumen   =======%
\keywords{Primera Palabra Clave \and Segunda Palabra Clave \and Otra.}


\section{Introducción}
La introducción de un artículo científico es una sección que proporciona una visión general del tema del artículo y justifica la realización de la investigación. Debe ser concisa, clara y bien organizada, y debe responder a las siguientes preguntas:\\
\begin{itemize}
    \item ¿Cuál es el tema del artículo?
    \item ¿Cuál es la importancia del tema?
    \item ¿Cuáles son los principales hallazgos del artículo?
\end{itemize}

La introducción suele dividirse en tres partes:\\
\begin{enumerate}
    \item Introducción general: Esta parte proporciona una descripción general del tema del artículo. Debe ser lo suficientemente amplia para que un lector interesado pueda entender el contexto del artículo, pero no debe ser tan amplia que se convierta en una revisión de la literatura.
Justificación de la investigación: Esta parte explica por qué el tema es importante y por qué es necesario realizar la investigación. Debe identificar los vacíos en el conocimiento o los desafíos que la investigación pretende abordar.
   \item Principales hallazgos: Esta parte proporciona una breve descripción de los principales hallazgos del artículo. No debe proporcionar detalles específicos, ya que estos se presentarán en las secciones siguientes del artículo.
   \item La introducción es una parte esencial de un artículo científico, ya que ayuda a los lectores a comprender el propósito y el alcance del artículo. Una introducción bien escrita puede ayudar a los lectores a decidir si el artículo es relevante para sus intereses y si vale la pena leerlo completo.
\end{enumerate}

A continuación se presentan algunos consejos para escribir una introducción de un artículo científico: \\

Empieza por identificar el tema principal del artículo. ¿Cuál es la pregunta o el problema que el artículo pretende abordar?\\
Una vez que hayas identificado el tema principal, explica por qué es importante. ¿Por qué es un tema relevante para la ciencia? ¿Por qué es importante para la sociedad?\\
Justifica la necesidad de realizar la investigación. ¿Cuáles son los vacíos en el conocimiento o los desafíos que la investigación pretende abordar?\\
Proporciona una breve descripción de los principales hallazgos del artículo. ¿Qué encontró la investigación?\\
Escribe la introducción de manera concisa y clara. La introducción debe ser fácil de entender y no debe contener información irrelevante.
\newline
Siguiendo estos consejos, podrás escribir una introducción de un artículo científico que sea concisa, clara y bien organizada.


\section{Encabezado sección.}
Bla bla bla de la Seccion.

\subsection{Encabezado de una subsección.}
Bla bla bla de la subseccion.
\newline
Y si escribimos una formulita?
\begin{equation}
F = \frac{G m M}{r^2}
\end{equation}
Para escribir ecuaciones en fomato \LaTeX  nos podemos ayudar de alguna GEN-IA como Chat-GPT \url{https://chat.openai.com/chat} o Bard \url{https://bard.google.com/}. Y tambien algo muy util es usar la siguiente pagina web. 
\url{https://latex.codecogs.com/eqneditor/editor.php?lang=es-es}


\section{Ejemplos de citas y figuras.}
\label{sec:others}

\subsection{Citas}
Si queremos poner algo explicito en el cuerpo del documento podemos escribir algo como:\\
Podemos encontrar la documentación o información en el siguiente link,
\begin{center}	\url{https://onepiece.fandom.com/es/wiki/One_Piece_Wiki}
\end{center}

O bien para citar uno tiene que escribir la cita en el archivo que se llama references.bib y luego en este documento el template.tex podemos citar de la siguiente manera:\\
\newline
Bla bla ... bueno el teorema de Pitagoras \cite{wikipedia-pitagoras} es feo.

\subsection{Figuras}
Ahora veremos como poner una figura: \\
bla bla miren la figura \ref{fig:figura1}
vieron la figura? es muy bonita.\\
Y ahora creamos la figura

\begin{figure}
	\centering
        \includegraphics[scale=0.8]{one_piece.png}
	\caption{Explicación de la figura.}
	\label{fig:figura1} % este es el nombre la figura
\end{figure}




\bibliographystyle{unsrtnat}
\bibliography{references} 













%%% Uncomment this line and comment out the ``thebibliography'' section below to use the external .bib file (using bibtex) .


%%% Uncomment this section and comment out the \bibliography{references} line above to use inline references.
% \begin{thebibliography}{1}

% 	\bibitem{kour2014real}
% 	George Kour and Raid Saabne.
% 	\newblock Real-time segmentation of on-line handwritten arabic script.
% 	\newblock In {\em Frontiers in Handwriting Recognition (ICFHR), 2014 14th
% 			International Conference on}, pages 417--422. IEEE, 2014.

% 	\bibitem{kour2014fast}
% 	George Kour and Raid Saabne.
% 	\newblock Fast classification of handwritten on-line arabic characters.
% 	\newblock In {\em Soft Computing and Pattern Recognition (SoCPaR), 2014 6th
% 			International Conference of}, pages 312--318. IEEE, 2014.

% 	\bibitem{hadash2018estimate}
% 	Guy Hadash, Einat Kermany, Boaz Carmeli, Ofer Lavi, George Kour, and Alon
% 	Jacovi.
% 	\newblock Estimate and replace: A novel approach to integrating deep neural
% 	networks with existing applications.
% 	\newblock {\em arXiv preprint arXiv:1804.09028}, 2018.

% \end{thebibliography}


\end{document}
